\begin{note}
    Если $g : E \to \overline{\R}$ измерима, то $\widetilde{g} : E \times \R \to \overline{\R}$, $\widetilde{g} (x, t) = g(x)$ также измерима. Это вытекает из того, что если
    $A = \{x \in E : g(x) < a\}$ измеримо, то и $\{(x, t) : \widetilde{g}(x, t) < a\} = A \times \R$ измеримо, что предлагается доказать в следующей задаче.
\end{note}

\begin{problem}
    Пусть $A$ измеримо, докажите измеримость $A \times \R$ (Указание: $A = \Delta \cup Z$).
\end{problem}

Рассмотрим одно важнейшее приложение принципа Кавальери.

\begin{corollary}
    Если $E \subset \R^n$ и $f : E \to [0, +\infty]$ измерима, то ее подграфик $S_f = \{(x, t) : x \in E, \ 0 < t < f(x)\}$ измерим в $\R^{n+1}$ и $\mu_{n+1}(S_f) = \int_{E} f \, d\mu_n$.
\end{corollary}

\begin{myproof}
    Функция $F(x, t) = t - f(x)$ измерима на $E \times \R$ как разность непрерывной функции $G(x,t) = t$ и $H(x, t) = f(x)$, измеримой по замечанию. 
    Тогда $S_f = \{(x, t) \in E \times \R : t > 0\} \cap \{(x, t) \in E \times \R : F(x, t) < 0\}$ измеримо в $\R^{n+1}$ как пересечение.
    Имеем \[(S_f)_x = \begin{cases}
     (0, f(x)), & x \in E, \\  
     \emptyset, & x \notin E. 
    \end{cases}\]
    Тогда $\forall x \in E : \mu_1 ((S_f)_x) = f(x)$, поэтому $\int_{E} f \, d\mu = \mu_{n+1} (S_f)$ по принципу Кавальери. 
\end{myproof}

\begin{problem}
    Докажите, что $\mu_{n+1} (S_f) = \int_{0}^{+\infty} \mu \{x \in E : f(x) > t\} \, dt$.
\end{problem}

\begin{example}
    Вычислим меру $n$-мерного шара $B_R(a)$. 
\end{example}

\begin{myproof}
    По свойству преобразования меры при гомотетии и сдвиге
    $\mu_n(B_R(a)) = \mu_n(B_R(0)) = R^n \mu_n(B_1(0))$. 
    Положим за $\omega_n = \mu_n(B_1(0))$ объем единичного шара в $\R^n$ и обозначим $B = B_1(0)$. 
    Если записать $\R^n$ в виде $\R^2_{(x_1, x_2)} \times \R^{n-2}_y$, то сечение $B_{(x_1, x_2)} = \{y \in \R^{n-2} : |y|^2 < 1 - x_1^2 - x_2^2\}$, то есть 
    \[ B_{(x_1, x_2)} = \begin{cases}
     B_{\sqrt{1 - x_1^2 - x_2^2}}(0), & x_1^2 + x_2^2 < 1,  \\  
     \emptyset, & x_1^2 +x_2^2 \geq 1. 
    \end{cases}\]
    Следовательно, по принципу Кавальери 
    \[ \omega_n = \int_{x_1^2 + x_2^2 < 1} \mu_{n-2} (B_{\sqrt{1-x_1^2-x_2^2}} (0)) \, dx_1 \, dx_2 = \omega_{n-2} \int_{x_1^2+x_2^2 < 1} (1 - x_1^2 - x_2^2)^{\frac{n-2}{2}} \, dx_1 \, dx_2. \] 
	Переходя к полярным координатам $(r, \phi)$, где $x_1^2+x_2^2 = r^2$ и $dx_1 \, dx_2 = r \, dr \, d\phi$ (замена в кратном интеграле производится по теореме 1.4 ниже), получаем:
    \[ \omega_n = \omega_{n-2} \int_0^{2\pi} d\phi \int_0^1 (1-r^2)^{\frac{n-2}{2}} r \, dr = 2\pi \omega_{n-2} \int_0^1 (1-r^2)^{\frac{n-2}{2}} r \, dr. \]
    Вычислим интеграл заменой $u = 1-r^2$, $du = -2rdr$:
    \[ \int_0^1 (1-r^2)^{\frac{n-2}{2}} r \, dr = \int_1^0 u^{\frac{n-2}{2}} \left(-\frac{1}{2} du\right) = \frac{1}{2} \int_0^1 u^{\frac{n-2}{2}} du = \frac{1}{2} \left[ \frac{u^{n/2}}{n/2} \right]_0^1 = \frac{1}{n}. \]
    Таким образом, мы получили рекуррентное соотношение $\omega_n = \frac{2\pi}{n} \omega_{n-2}$.
    Используя базовые случаи $\omega_1=2$ и $\omega_2=\pi$ можем найти объём для любой размерности $n$.
\end{myproof}


\begin{problem}
    Пусть $T : \R^n \to \R^n$ обратимый линейный оператор и $E \subset \R^n$ измеримо. 
    Докажите, что $\mu_n(T(E)) = |\det T| \cdot \mu_n(E)$ (Указание: \href{https://youtu.be/MrbVp30tSPc?list=PL4_hYwCyhAvaNULcW2kwdthjEfCDAMsRy&t=3172}{доказательство Редкозубова}).
\end{problem}

\subsection{Замена переменных в интеграле}

\begin{definition}
    Пусть $U$ и $V$ открыты в $\R^n$. Отображение $g : U \to V$ называется \textit{диффеоморфизмом}, если $g$ -- биекция, $g \in C^1(U, V)$, и $g^{-1} \in C^1(V, U)$.
\end{definition}

Всюду в дальнейшем под $Dg(x)$ понимается матрица Якоби гладкого отображения $g$ в точке $x$. 
Определитель $J_g(x) = \det Dg(x)$ называется \textit{якобианом} отображения $g$ в точке $x$.

\begin{note}
    Если $g$ -- диффеоморфизм, то $g^{-1} \circ g = \id_U$. По теореме о дифференцировании композиции $Dg^{-1}(g(x)) \cdot Dg(x) = E$. 
    Следовательно, матрица $Dg(x)$ невырождена и $Dg^{-1}(g(x)) = [ Dg(x) ]^{-1}$.
\end{note}

\begin{theorem}
    Пусть $g : U \to V$ -- диффеоморфизм открытых $U, V \subset \R^n$ и $E \subset U$ измеримо. 
    Если $f$ неотрицательная измеримая или интегрируемая на $g(E)$, то 
    \[ \int_{g(E)} f(y) \, d\mu(y) = \int_{E} (f \circ g)(x) \cdot |\det Dg(x)| \, d\mu(x). \]
\end{theorem}

Сначала докажем несколько технических результатов. 

\begin{lemma}
    Пусть $x_0 \in U$ и $\epsilon > 0$. Тогда найдется $\delta > 0$, что для всякого замкнутого куба $Q$, $x_0 \in Q \subset U$ с длиной ребра $l(Q) < \delta$, справедливо неравенство 
    \[ \frac{\mu(g(Q))}{\mu(Q)} \leq |J_g(x_0)| + \epsilon. \] 
\end{lemma}

\begin{myproof}
    Сначала отметим, что $g(Q)$ измеримо, так как является образом компакта при непрерывном отображении. 
    Также будем использовать норму $||x|| = \max \limits_{1 \leq i \leq n} |x_i|$ на $\R^n$, так как все нормы эквивалентны. 
    По этой норме <<шар>> $C_r(y) = \{x : ||x - y|| \leq r\}$ будет замкнутым кубом с длиной ребра $l(C) = 2r$.

    \vspace{5pt}
    Сначала рассмотрим случай $Dg(x_0) = E$ и $g(x_0) = 0$.
	В этом случае $|J_g(x_0)| = 1$.
    По определению дифференцируемости в точке $x_0$ имеем 
	\[ g(x) = g(x_0) + Dg(x_0)(x-x_0) + \alpha(x) \cdot ||x - x_0|| = x-x_0 + \alpha(x) \cdot ||x-x_0||, \ \alpha(x) \to 0 \text{ при } x \to x_0. \]
    Заранее выберем $\eta > 0$ так, что $(1+2\eta)^n < 1 + \epsilon$ и пусть $\delta > 0$ таково, что если $||x - x_0|| < \delta$, то $||\alpha(x)|| < \eta$.
    Теперь возьмем любой замкнутый куб $Q$, где $x_0 \in Q \subset U$ и $l(Q) =: s < \delta$. 
    Пусть $a$ -- центр $Q$. Тогда $Q = C_{s/2} (a)$. 
    Для любой точки $x \in Q$ по неравенству треугольника 
	\[ ||x - x_0|| \leq ||x - a|| + ||a - x_0|| \leq s/2 + s/2 = s < \delta. \]
	Раз $||x - x_0|| < \delta$, то $||\alpha(x)|| < \eta$. 
	Теперь оценим диаметр образа $g(Q)$. Возьмем $x_1, x_2 \in Q$.
	\[ || g(x_1) - g(x_2) || = || \, x_1 - x_2 + \alpha(x_1) \cdot ||x_1-x_0|| - \alpha(x_2) \cdot ||x_2-x_0|| \, ||. \]
	Подставляем оценки и используем неравенство треугольника: 
	\[ || g(x_1) - g(x_2) || \leq ||x_1 - x_2|| + ||\alpha(x_1)|| \cdot ||x_1-x_0|| + ||\alpha(x_2)|| \cdot ||x_2-x_0|| < s + \eta \cdot s + \eta \cdot s = s(1+2\eta).\]
	Следовательно, $g(Q)$ содержится в некотором кубе $Q'$ с $l(Q') = s(1 + 2\eta)$ \footnote{$g(Q) \subset Q' := [a_1, a_1 + s(1 + 2\eta)] \times \dotsc \times [a_n, a_n + s(1 + 2\eta)]$, где $a_i = \inf \{y_i : y \in g(Q)\}$, так как в противном случае найдутся $y', y'' \in g(Q)$ для которых неверно доказанное неравенство.}. 
	По монотонности меры $\mu(g(Q)) \leq \mu(Q') = s^n (1+2\eta)^n$. 
    Учитывая, что $\mu(Q) = s^n$, получаем
    \[ \frac{\mu(g(Q))}{\mu(Q)} \leq \frac{s^n(1+2\eta)^n}{s^n} \leq (1+2\eta)^n < 1 + \epsilon. \]
	Для простого случая лемма доказана.

    \vspace{5pt}
    Общий случай, пусть $T = Dg(x_0)$, $T$ невырождена. Рассмотрим $h = T^{-1} g - T^{-1} g(x_0)$ диффеоморфизм как композиция диффеоморфизмов.
	Проверим ее свойства в точке $x_0$:
	\begin{itemize}[itemsep = 0pt, topsep = 5pt]
		\item $h(x_0) = T^{-1}(g(x_0)) - T^{-1}(g(x_0)) = 0$.
		\item $Dh(x_0) = D(T^{-1} \circ g)(x_0) = T^{-1} \cdot Dg(x_0) = T^{-1} \cdot T = E$.
	\end{itemize}
    Следовательно, $h$ удовлетворяет условиям простого случая. Зададим погрешность $\epsilon' = \epsilon / |\det T|$, тогда $\exists \delta > 0$, что для любого куба $Q$ с $x_0 \in Q$ и $l(Q) < \delta$ 
	выполняется: 
    \[ \frac{\mu(h(Q))}{\mu(Q)} \leq |J_h(x_0)| + \epsilon' = 1 + \frac{\epsilon}{|\det T|}. \]
    По задаче о преобразовании меры линейным оператором и из инвариантности меры относительно сдвига имеем $\mu(h(Q)) = \mu(T^{-1} g(Q)) = |\det T^{-1}| \cdot \mu(g(Q)) = \frac{\mu(g(Q))}{|\det T|}$. 
    Подставляя в неравенство для $h$ и умножая на $|\det T| = |J_g(x_0)|$, получаем искомое.
 \end{myproof}


 \begin{corollary}
    Пусть $\{Q_k\}$ -- последовательность замкнутых вложенных кубов ($Q_k \supset Q_{k+1}$) и $l(Q_k) \to 0$. 
    Если $x_0$ -- единственная точка из $\bigcap Q_k$, то 
    \[ \varlimsup \limits_{k \to \infty} \frac{\mu(g(Q_k))}{\mu(Q_k)} \leq |J_g(x_0)|. \]
 \end{corollary}

 \begin{lemma}
    Если $G$ открыто в $U$, то $g(G)$ открыто в $V$ и 
    \[ \mu(g(G)) \leq \int_{G} |J_g(x)| \, d\mu(x). \]
 \end{lemma}

 \begin{myproof}
    Множество $g(G)$ открыто, так как $g(G) = (g^{-1})^{-1}(G)$. Пусть $G = \bigsqcup \limits_{k=1}^{\infty} B_k$ -- объединение двоичных кубов. 
	Из инъективности $g(G) = \bigsqcup g(B_k)$.
	По свойству счетной аддитивности меры Лебега и интеграла 
	\[ \mu(g(G)) = \sum_{k=1}^{\infty} \mu(g(B_k)), \quad  \int_G |J_g(x)| \, d\mu(x) = \sum_{k=1}^{\infty} \int_{B_k} |J_g(x)| \, d\mu(x). \]
	Поэтому докажем $\mu(g(B_k)) \le \int_{B_k} |J_g(x)| \, d\mu(x)$ для всех $k$ и просуммируем неравенства ($g(B_k)$ измеримы, так как $B_k$ представимы счетным объединением компактов).

    \vspace{5pt}
    Предположим, что лемма для куба неверна. Тогда найдется такой $Q$, что неравенство неверно и при этом $\overline{Q} \subset U$. Мы можем усилить это неравенство: $\exists \epsilon > 0$, такое что
    \[ \mu(g(Q)) > \int_{Q} |J_g(x)| \, d\mu(x) + \epsilon \mu(Q). \]
    Деля каждое ребро пополам, получим разбиение $Q$ на $2^n$ двоичных кубов $\{Q_i\}_{i=1}^{2^n}$. 
    Используя конечную аддитивность меры и интеграла заключаем, что хотя бы для одного из кубов $Q_i$ (обозначим его через $Q_1$) выполнено неравенство со знаком <<$>$>>. 
	Повторяя этот процесс, мы строим последовательность вложенных замкнутых кубов $C_k$, где $C_1 = \overline{Q}$, $C_2 = \overline{Q_1}$ и так далее, где $l(C_k) \to 0$ и $\bigcap C_k = \{x_0\}$.
    Тогда для всех $k$ 
    \[ \mu(g(C_k)) > \int_{C_k} |J_g(x)| \, d\mu(x) + \epsilon \mu(C_k). \]
    Делим на $\mu(C_k)$. Так как $J_g$ -- непрерывная функция (якобиан является полиномиальной комбинацией частных производных, непрерывных по определению диффеоморфизма), то $\frac{1}{\mu(C_k)} \cdot \int_{C_k} |J_g(x)| \, d\mu(x) \underset{k \to \infty}{\to} |J_g(x_0)|$ (из определения непрерывности в $x_0$), откуда $\varlimsup \limits_{k \to \infty} \frac{\mu(g(C_k))}{\mu(C_k)} \geq |J_g(x_0)| + \epsilon$.
    Противоречие с предыдущим следствием.
 \end{myproof}
